\documentclass[11pt]{article}
\usepackage[a4paper,margin=23mm]{geometry}
\usepackage{amsmath,amssymb,amsthm,mathtools,bm}
\usepackage{booktabs,longtable,array}
\usepackage{enumitem}
\usepackage{xcolor,microtype,xurl}
\usepackage[hidelinks]{hyperref}
\setlength{\emergencystretch}{3em}
\numberwithin{equation}{section}
\newtheorem{theorem}{Theorem}[section]
\newtheorem{proposition}[theorem]{Proposition}
\newtheorem{lemma}[theorem]{Lemma}
\newtheorem{corollary}[theorem]{Corollary}
\theoremstyle{definition}
\newtheorem{definition}[theorem]{Definition}
\newtheorem{gate}[theorem]{Fail-closed gate}
\theoremstyle{remark}
\newtheorem{remark}[theorem]{Remark}
\newcommand{\RR}{\mathbb R}
\newcommand{\ZZ}{\mathbb Z}
\newcommand{\cA}{\mathcal A}
\newcommand{\cC}{\mathcal C}
\newcommand{\cE}{\mathcal E}
\newcommand{\dd}{\mathrm d}
\newcommand{\Tr}{\operatorname{Tr}}
\newcommand{\diag}{\operatorname{diag}}

\title{AP-E8 Topology-Preserving Lattice Action:\\
An Unanchored Finite-Grid \(B=1\) Stationary Family\\
with a Fail-Closed Continuum Boundary}
\author{Route F research note}
\date{17 July 2026}

\begin{document}
\maketitle

\begin{abstract}
AP-E7 proved that the unrestricted fixed-boundary finite-site space is
connected, so the naive sampled action has no exact integer topological
sectors and its unanchored relaxations unwind.  This note changes that
specific blocker.  On a fixed Freudenthal triangulation we add a logarithmic
barrier to every unique pair of vertices sharing a tetrahedron.  Its linear
term is subtracted, so it diverges at a pairwise admissibility floor that
prevents reaching the dislocation locus while contributing
only \(O(a^2)\) to a uniformly smooth sampled field.  Pairwise admissibility
gives an explicit \(0.50744\) lower bound on the norm of every affine
tetrahedron.  The normalized interpolation is therefore nonsingular, its
degree is locally constant, and no continuous finite-energy path changes
\(B\).

For every fixed grid and every nonempty admissible component, compactness of
the sphere variables, coercivity of the breathing field, and divergence of
the barrier prove existence of an interior unanchored minimizer.  Sampling a
smooth continuum degree-one field shows that such \(B=1\) components exist on
all sufficiently fine grids.  Numerically, seven independently re-relaxed
representatives at four spacings and three volumes have direct full tangent
gradient densities below \(8\times10^{-7}\), three-target charge
\((1,1,1)\), and positive barrier margins without a core or centre pin.
Nearby barrier parameters give the same finite-grid conclusion.

The result is deliberately narrower than a continuum theorem.  The
grid-indexed minimizers have not been proved equicoercive, convergent, barrier-
independent, or triangulation-independent.  No physical two-colour gauge
measure, determinant Hessian, BRST superdeterminant, HMC calculation, or
degree-one Route-E portal is supplied.  The finite-grid topology/stationarity
sub-blocker is closed; the physical Lane-B conjunction remains open.
\end{abstract}

\tableofcontents

\section{Decision and logical scope}

The three proposed blocker-changing paths were:
\begin{enumerate}[label=(\roman*)]
\item complete a finite Ginsparg--Wilson/domain-wall target kernel;
\item replace the naive lattice action by a topology-preserving one and find
an unanchored stationary sequence;
\item compute the actual \(SO(3)\) Yukawa mapping-torus mod-two index.
\end{enumerate}
The first route still lacks one target mass-family lift, its mixed determinant
orientation, and an all-scale \(k=+2\) mechanism.  The third route can select a
Pfaffian torsion line but does not repair the finite-site unwinding found in
AP-E7.  Route (ii) is selected because it directly changes the proved
finite-site obstruction: the infinite barrier removes the dislocation locus
from the finite-energy configuration space and creates separated components.

This construction is inspired by geometric lattice charge and lattice
Skyrmion admissibility ideas \cite{BergLuscher1981APE8,
Berg1981DislocationsAPE8,Ward1995APE8,SchrammSvetitsky2000APE8}, but the
specific action below uses a subtracted logarithm and counts every unordered
Freudenthal co-cell pair only once.  Its finite-energy boundary lies strictly
inside the normalized-affine admissible region, before any true dislocation.
No claim of regulator uniqueness is
made.  Recent dense two-colour effective theories show that diquark, scalar,
vector, and axial channels can all affect the dynamics
\cite{SuenagaEtAl2023APE8,SuenagaEtAl2024APE8}; they do not convert the present
three-dimensional classical regulator into four-dimensional QC2D.

\section{Finite-grid data and the AP-E7 base action}

Let \(\Lambda_{a,L}\) be an \(N^3\) cubic grid with spacing
\(a=L/(N-1)\).  The outer faces are fixed to
\begin{equation}
 n_x=e_0=(1,0,0,0),\qquad s_x=1,
 \qquad x\in\partial\Lambda_{a,L}.
 \label{eq:boundary}
\end{equation}
At each of the \(M=(N-2)^3\) interior sites,
\(n_x\in S^3\subset\RR^4\) and \(s_x\in\RR\).  Before admissibility, the
configuration space is
\begin{equation}
 \cC_N=(S^3)^M\times\RR^M,
 \label{eq:naive-space}
\end{equation}
which is path-connected.

The AP-E7 base action is retained literally:
\begin{align}
 E_a[n,s]={}&\frac a2\sum_{\langle xy\rangle}
 \bigl(\lVert n_y-n_x\rVert^2+(s_y-s_x)^2\bigr)\nonumber\\
 &+\frac{\kappa a^3}{2}\sum_c\sum_{i<j}
 \left(\lVert u_i^c\rVert^2\lVert u_j^c\rVert^2
 -(u_i^c\!\cdot u_j^c)^2\right)\nonumber\\
 &+a^3\sum_{x\in\mathrm{int}}
 \left[\frac{m_s^2}{2}(s_x-1)^2
 +\beta^2s_x^2(1-n_{x0})\right],
 \label{eq:base-action}
\end{align}
where
\begin{equation}
 u_i^c=\frac{n_{c+\hat\imath}-n_c}{a},\qquad
 \kappa=1,\quad\beta=\frac12,\quad m_s=\frac52.
 \label{eq:base-parameters}
\end{equation}
Every term in Eq.~\eqref{eq:base-action} is nonnegative.  The last line is
coercive in all interior \(s_x\).  This is a classical Skyrme-type regulator
\cite{Skyrme1961APE8}, not a lattice gauge-theory action.

Each cube is split into the six Freudenthal tetrahedra.  For a permutation
\(p\in S_3\), their vertices are
\begin{equation}
 x_0=c,\qquad x_1=c+e_{p_1},\qquad
 x_2=c+e_{p_1}+e_{p_2},\qquad x_3=c+(1,1,1).
 \label{eq:freudenthal}
\end{equation}
Let \(\cE_F\) be the set of \emph{unique unordered} vertex pairs
\(\{x,y\}\) that co-occur in at least one such tetrahedron.  Pair incidence
is not multiplied when the same edge belongs to several tetrahedra.

\section{The topology-preserving action}

Fix
\begin{equation}
 -\frac13<\epsilon<1,\qquad \gamma>0,\qquad
 q_{xy}=n_x\cdot n_y.
\end{equation}
For the production card, \(\epsilon=0.01\) and \(\gamma=1\).  Put
\begin{equation}
 b_\epsilon(t)=
 \begin{cases}
 -\log\!\dfrac{t-\epsilon}{1-\epsilon}
 +\dfrac{t-1}{1-\epsilon},&t>\epsilon,\\[6pt]
 +\infty,&t\le\epsilon.
 \end{cases}
 \label{eq:barrier}
\end{equation}
The AP-E8 action is
\begin{equation}
 \boxed{
 E^{\mathrm{TP}}_{a,\epsilon,\gamma}[n,s]
 =E_a[n,s]+\gamma a\sum_{\{x,y\}\in\cE_F}
 b_\epsilon(q_{xy}).}
 \label{eq:tp-action}
\end{equation}

The subtraction in Eq.~\eqref{eq:barrier} is essential.  With
\(d=1-\epsilon\), direct differentiation gives
\begin{align}
 b_\epsilon(1)&=b_\epsilon'(1)=0,
 &b_\epsilon'(t)&=-\frac1{t-\epsilon}+\frac1d
 =\frac{t-1}{d(t-\epsilon)},\label{eq:bprime}\\
 b_\epsilon''(t)&=\frac1{(t-\epsilon)^2}>0,
 &\lim_{t\downarrow\epsilon}b_\epsilon(t)&=+\infty.
 \label{eq:bsecond}
\end{align}
Since \(t\le1\) for unit vectors, \(b_\epsilon\) decreases from infinity to
zero and is nonnegative on its finite domain.  Writing
\(u=1-t\in[0,d)\),
\begin{equation}
 b_\epsilon(1-u)
 =-\log\left(1-\frac ud\right)-\frac ud
 =\sum_{k=2}^{\infty}\frac1k\left(\frac ud\right)^k
 =\frac{u^2}{2d^2}+O(u^3).
 \label{eq:barrier-series}
\end{equation}
Thus no term linear in \(1-t\) remains.

\section{No-zero lemma and exact finite-energy sectors}

\begin{lemma}[Pairwise admissibility implies a uniform no-zero bound]
Let \(n_0,\ldots,n_3\in S^3\) obey
\(n_i\cdot n_j>\epsilon\) for every \(i\ne j\).  For all barycentric
coordinates \(\lambda_i\ge0\), \(\sum_i\lambda_i=1\),
\begin{equation}
 \left\lVert\sum_{i=0}^3\lambda_i n_i\right\rVert^2
 \ge \epsilon+(1-\epsilon)\sum_{i=0}^3\lambda_i^2
 \ge\frac{1+3\epsilon}{4}>0.
 \label{eq:no-zero-bound}
\end{equation}
\end{lemma}

\begin{proof}
Using \(2\sum_{i<j}\lambda_i\lambda_j
=1-\sum_i\lambda_i^2\),
\begin{align}
 \left\lVert\sum_i\lambda_i n_i\right\rVert^2
 &=\sum_i\lambda_i^2
 +2\sum_{i<j}\lambda_i\lambda_j(n_i\cdot n_j)\nonumber\\
 &\ge\sum_i\lambda_i^2
 +2\epsilon\sum_{i<j}\lambda_i\lambda_j
 =\epsilon+(1-\epsilon)\sum_i\lambda_i^2.
\end{align}
The first inequality is strict whenever at least two barycentric weights are
positive; at a vertex both sides of that comparison equal one.  Finally,
the Cauchy--Schwarz inequality gives \(\sum_i\lambda_i^2\ge1/4\).
The last bound is positive exactly when \(\epsilon>-1/3\).
\end{proof}

For \(\epsilon=0.01\), Eq.~\eqref{eq:no-zero-bound} yields
\begin{equation}
 \min_T\min_{\lambda\in\Delta^3}
 \left\lVert\sum_i\lambda_i n_i\right\rVert
 \ge\sqrt{\frac{1+3\epsilon}{4}}
 =0.507444578255\ldots.
 \label{eq:numerical-lower-bound}
\end{equation}
Therefore the normalized affine interpolation
\begin{equation}
 \widehat n_T(\lambda)
 =\frac{\sum_i\lambda_i n_i}
 {\lVert\sum_i\lambda_i n_i\rVert}
 \label{eq:normalized-affine}
\end{equation}
is continuous on every tetrahedron, including common faces.

Define the admissible open set
\begin{equation}
 \cA_{\epsilon,N}
 =\{(n,s)\in\cC_N:q_{xy}>\epsilon
 \text{ for all }\{x,y\}\in\cE_F\}.
 \label{eq:admissible-set}
\end{equation}
Collapsing the vacuum boundary to a point turns Eq.~\eqref{eq:normalized-affine}
into a continuous map \(S^3\to S^3\).  Its degree
\(B_{\rm geom}\in\ZZ\) is constant on every connected component of
\(\cA_{\epsilon,N}\).

\begin{theorem}[Finite-energy sector preservation]
Every continuous path with finite
\(E^{\mathrm{TP}}_{a,\epsilon,\gamma}\) remains in one component of
\(\cA_{\epsilon,N}\), and its normalized-affine degree is constant.
Consequently no continuous finite-energy path joins \(B=1\) to the vacuum.
\end{theorem}

\begin{proof}
Finite action requires every \(q_{xy}>\epsilon\), so the entire path lies in
\(\cA_{\epsilon,N}\).  A connected path cannot change connected component.
Equation~\eqref{eq:no-zero-bound} makes the normalized interpolation a
continuous homotopy along the path, and homotopy invariance fixes its degree.
Equivalently, any route to another degree must meet some
\(q_{xy}=\epsilon\), where Eq.~\eqref{eq:barrier} has infinite energy.
\end{proof}

\section{Existence of an unanchored finite-grid stationary family}

\begin{theorem}[Fixed-grid minimizer]
For fixed \(a>0\), fixed finite \(N\), \(\gamma>0\), \(m_s>0\), and any
nonempty connected component \(K\subset\cA_{\epsilon,N}\), the action
\(E^{\mathrm{TP}}_{a,\epsilon,\gamma}\) attains a minimum at an interior
point of \(K\).  That minimizer is stationary under every interior
\(S^3\)-tangent variation of \(n\) and every interior variation of \(s\).
No core or centre anchor is required.
\end{theorem}

\begin{proof}
Choose a minimizing sequence in \(K\).  After discarding finitely many terms,
its energy is bounded by a finite constant \(C\).  Nonnegativity of every
term in Eq.~\eqref{eq:tp-action} implies, for every edge,
\begin{equation}
 b_\epsilon(q_{xy})\le\frac{C}{\gamma a}.
 \label{eq:barrier-sublevel}
\end{equation}
The strictly decreasing map
\(b_\epsilon:(\epsilon,1]\to[0,\infty)\) is invertible.  Hence there is a
fixed-grid \(\delta_{a,C}>0\) such that
\begin{equation}
 q_{xy}\ge\epsilon+\delta_{a,C}
 \label{eq:sublevel-margin}
\end{equation}
for every member of the sequence and every edge.  All \(n_x\) already lie in
the compact finite product \((S^3)^M\).  The mass term gives
\begin{equation}
 \frac{a^3m_s^2}{2}\sum_{x\in\mathrm{int}}(s_x-1)^2\le C,
 \label{eq:s-coercive}
\end{equation}
so the finite vector \((s_x)\) is bounded.  A subsequence converges.  The
uniform margin \eqref{eq:sublevel-margin} keeps its limit inside
\(\cA_{\epsilon,N}\); connected components are closed relative to that open
set, so the limit remains in \(K\).  Continuity on the finite-energy domain
shows that the limit attains the infimum.

The margin also supplies a neighborhood contained in the same admissible
component.  The minimizer is therefore an unconstrained interior critical
point of the fixed-boundary finite-dimensional manifold.  Its tangent and
scalar first variations vanish.  No value at an interior site was fixed in
the argument.
\end{proof}

\begin{theorem}[Nonempty \(B=1\) components on a refining grid sequence]
Let \(n:\overline\Omega\to S^3\) be a smooth degree-one map equal to the
vacuum near \(\partial\Omega\), and let \(s\) be smooth.  For any fixed
\(\epsilon<1\), all sufficiently fine Freudenthal grids obtained by sampling
\((n,s)\) belong to \(\cA_{\epsilon,N}\) and their normalized-affine degree
is one.  Consequently Theorem 5.1 supplies an unanchored stationary minimizer
\((n_a,s_a)\) in a \(B=1\) component for every grid in a sequence
\(a\downarrow0\).
\end{theorem}

\begin{proof}
Uniform continuity of \(n\) on the compact domain implies that the images of
any two vertices sharing a tetrahedron approach one another uniformly as
\(a\to0\).  Their dot products therefore approach one uniformly and exceed
\(\epsilon\) for sufficiently small \(a\).  The piecewise affine
interpolation of the sampled values converges uniformly to \(n\), and its
norm converges uniformly to one.  After normalization it is uniformly close
to \(n\); for sufficiently small \(a\), the pointwise shortest-geodesic
homotopy is well-defined.  Thus its degree equals \(\deg n=1\).  Apply the
fixed-grid theorem to the component containing that sample.
\end{proof}

\begin{remark}[What ``sequence'' does and does not mean]
The theorem gives a family of exact finite-dimensional minimizers indexed by
the grid spacing.  It does not give a compactness estimate uniform in \(a\).
Indeed, Eq.~\eqref{eq:barrier-sublevel} permits
\(\delta_{a,C}\) to shrink rapidly as \(a\downarrow0\).  The theorem alone
does not imply that \((n_a,s_a)\) converges, that its limit solves a continuum
Euler--Lagrange equation, or that a barrier-supported lattice core cannot
form.
\end{remark}

\section{Smooth-field scaling of the new term}

Let an edge have physical displacement \(av\) with \(v\) fixed by the
Freudenthal stencil.  For a uniformly smooth unit field,
\begin{equation}
 n(x)\cdot n(x+av)
 =1-\frac{a^2}{2}\lVert D_vn(x)\rVert^2+O(a^3).
 \label{eq:dot-expansion}
\end{equation}
Combining Eqs.~\eqref{eq:barrier-series} and \eqref{eq:dot-expansion},
\begin{equation}
 b_\epsilon(n(x)\cdot n(x+av))
 =\frac{a^4}{8(1-\epsilon)^2}
 \lVert D_vn(x)\rVert^4+O(a^5).
 \label{eq:edge-expansion}
\end{equation}
There are \(O(a^{-3})\) unique edges at fixed physical volume.  Therefore
\begin{equation}
 \gamma a\sum_{\{x,y\}\in\cE_F}b_\epsilon(q_{xy})
 =\frac{\gamma a^2}{8(1-\epsilon)^2}
 \int_\Omega\sum_v\rho_v\lVert D_vn\rVert^4\,\dd^3x
 +O(a^3),
 \label{eq:smooth-irrelevance}
\end{equation}
where \(\rho_v\) are the direction/incidence densities of the chosen unique-
edge stencil.  Thus the new term has no leading kinetic renormalization and
is \(O(a^2)\) on a fixed uniformly smooth family.

Equation~\eqref{eq:smooth-irrelevance} is not uniform over the unknown
minimizers.  It proves neither equicoercivity nor Gamma convergence.  It also
does not remove finite-\(a\) triangulation dependence, because the
\(\rho_v\) depend on the chosen stencil.

\section{Exact first and second variations}

Let \(w=\gamma a\), \(q=n_x\cdot n_y\), and consider one barrier edge
\(w b_\epsilon(q)\).  Its ambient gradients are
\begin{equation}
 G_x^{(xy)}=w b_\epsilon'(q)n_y,\qquad
 G_y^{(xy)}=w b_\epsilon'(q)n_x.
 \label{eq:edge-gradient}
\end{equation}
The ambient Hessian blocks are
\begin{align}
 H_{xx}^{(xy)}&=w b_\epsilon''(q)n_yn_y^T,&
 H_{yy}^{(xy)}&=w b_\epsilon''(q)n_xn_x^T,\nonumber\\
 H_{xy}^{(xy)}&=w\left[b_\epsilon''(q)n_yn_x^T
 +b_\epsilon'(q)I_4\right],&
 H_{yx}^{(xy)}&=(H_{xy}^{(xy)})^T.
 \label{eq:edge-ambient-hessian}
\end{align}

For tangent vectors \(\xi_x\perp n_x\), \(\xi_y\perp n_y\), take the
product geodesic on \(S^3\times S^3\).  Its first dot-product variation is
\begin{equation}
 \delta q=\xi_x\cdot n_y+n_x\cdot\xi_y,
 \label{eq:delta-q}
\end{equation}
and its second variation is
\begin{equation}
 \delta^2q=2\xi_x\cdot\xi_y
 -q\bigl(\lVert\xi_x\rVert^2+\lVert\xi_y\rVert^2\bigr).
 \label{eq:delta2-q}
\end{equation}
Hence the exact intrinsic edge quadratic form is
\begin{align}
 \delta^2\bigl(wb_\epsilon(q)\bigr)
 =w\Bigl\{&b_\epsilon''(q)
 (\xi_x\cdot n_y+n_x\cdot\xi_y)^2\nonumber\\
 &+b_\epsilon'(q)
 \left[2\xi_x\cdot\xi_y
 -q(\lVert\xi_x\rVert^2+\lVert\xi_y\rVert^2)\right]\Bigr\}.
 \label{eq:intrinsic-edge-hessian}
\end{align}

For the full action, let \(G_x=\partial E^{\rm TP}/\partial n_x\),
\(P_x=I_4-n_xn_x^T\), and \(\lambda_x=n_x\cdot G_x\).  Stationarity is
\begin{equation}
 P_xG_x=0,\qquad \frac{\partial E^{\rm TP}}{\partial s_x}=0.
 \label{eq:stationarity}
\end{equation}
If \(Q\) collects orthonormal tangent frames and the scalar directions, the
full Riemann Hessian at a stationary point is
\begin{equation}
 H_{\rm Riem}
 =Q^T\left[H_{\rm emb}
 -\diag(\lambda_1I_4,0,\lambda_2I_4,0,\ldots)\right]Q.
 \label{eq:riemann-hessian}
\end{equation}
Projecting \(H_{\rm emb}\) without the sphere-curvature shift is incorrect.
The AP-E8 card validates Eq.~\eqref{eq:intrinsic-edge-hessian} independently;
it does not assemble or promote Eq.~\eqref{eq:riemann-hessian} as a physical
gauge--meson--ghost--fermion super-Hessian.

\section{Numerical protocol}

The reproducible implementation is
\begin{center}
\path{route_f/code/scan_ap_e8_topology_preserving_b1.py}.
\end{center}
It imports the canonical AP-E6 continuum profile only as initial data.  The
little-endian \((r,F,s)\) checksum is
\begin{center}
\small\texttt{81104f59a5f3f4337739fc2a217cafc8da91581c9f1fb40b4fdba6ce0f948d59}.
\end{center}
Every grid receives exact vacuum boundary data and is minimized independently.
No core, centre, or interior field value is fixed.

A fixed-base exponential chart is used within each L-BFGS solve and is
re-centred if needed.  An invalid trial point receives a solver-only
continuation which matches the exact \(b,b',b''\) at a positive cutoff and
adds a differentiable quartic rejection below the admissibility floor.  Its
energy/gradient pair is checked separately.  All reported fields are
re-evaluated against the literal infinite-barrier domain.  Because the
complete geodesic segment of
each internal quasi-Newton line search was not interval-certified,
\texttt{complete\_line\_search\_segments\_certified\_admissible=false}.
This does not affect the final-point stationarity test or the existence
theorem, but it forbids a stronger claim about the numerical trajectory.

At every final field the code discards the chart derivative and recomputes
the direct projected gradient
\begin{equation}
 g_a=\frac{
 \left[\sum_x\lVert P_xG_x\rVert^2
 +\sum_x\left(\partial E^{\rm TP}/\partial s_x\right)^2\right]^{1/2}}
 {a^3\sqrt{4M}}.
 \label{eq:gradient-density}
\end{equation}
The stationarity threshold is \(g_a<2\times10^{-6}\).  The degree is counted
independently with three generic regular targets.  A size diagnostic is
\begin{equation}
 R_{\rm rms}^2
 =\frac{\sum_x(1-n_{x0})|x|^2}{\sum_x(1-n_{x0})}.
 \label{eq:rms-radius}
\end{equation}
It is a profile diagnostic, not a proof of localization or continuum
compactness.

Three independent calculus regressions precede the relaxation:
\begin{align}
 \text{total-gradient directional residual}&=9.38\times10^{-11},\nonumber\\
 \text{invalid-continuation gradient residual}&=2.81\times10^{-11},\nonumber\\
 \text{barrier geodesic-Hessian residual}&=8.16\times10^{-9}.
 \label{eq:calculus-residuals}
\end{align}

\section{Numerical results}

\subsection{Four spacings and three volumes}

The seven principal solutions are:
\begin{center}
\small
\begin{tabular}{rrrrrrrr}
\toprule
\(N\)&\(L\)&\(a\)&\(E^{\rm TP}\)&\(g_a\)&\(\min q\)&\(B\) rows&\(R_{\rm rms}/a\)\\
\midrule
15&6&0.428571&75.090097547&\(1.651\!\times10^{-7}\)&0.091845&\((1,1,1)\)&3.0704\\
17&6&0.375000&72.444298275&\(3.421\!\times10^{-7}\)&0.090442&\((1,1,1)\)&3.3504\\
19&6&0.333333&70.831336466&\(2.028\!\times10^{-7}\)&0.083178&\((1,1,1)\)&3.6277\\
21&6&0.300000&69.844122881&\(7.704\!\times10^{-7}\)&0.077786&\((1,1,1)\)&3.8980\\
16&6&0.400000&73.627393329&\(2.150\!\times10^{-7}\)&0.090722&\((1,1,1)\)&3.2634\\
21&8&0.400000&73.346860060&\(1.052\!\times10^{-7}\)&0.095418&\((1,1,1)\)&3.5780\\
26&10&0.400000&73.289478682&\(1.228\!\times10^{-7}\)&0.095975&\((1,1,1)\)&3.8613\\
\bottomrule
\end{tabular}
\end{center}
All pair margins \(\min q-\epsilon\) are positive.  The fixed-
\(a=0.4\) energy spread over \(L=6,8,10\) is
\begin{equation}
 \frac{E_{\max}-E_{\min}}{E_{\max}}=0.00458952.
 \label{eq:volume-spread}
\end{equation}
The fixed-\(L=6\) physical radii for \(N=15,17,19,21\) are
\(1.3159,1.2564,1.2092,1.1694\), while \(R_{\rm rms}/a\) grows.  This is evidence
against a literal one-cell collapse in the displayed sequence, but it is not
a uniform lower-radius theorem.

\subsection{Barrier-parameter controls}

At \(N=15,L=6\), four nearby choices give:
\begin{center}
\small
\begin{tabular}{rrrrrr}
\toprule
\(\gamma\)&\(\epsilon\)&\(E^{\rm TP}\)&\(g_a\)&\(\min q\)&\(B\) rows\\
\midrule
0.5&0.010&63.665861079&\(1.614\!\times10^{-7}\)&0.044767&\((1,1,1)\)\\
2.0&0.010&90.696115305&\(1.794\!\times10^{-7}\)&0.221220&\((1,1,1)\)\\
1.0&0.005&74.773915108&\(1.487\!\times10^{-7}\)&0.085519&\((1,1,1)\)\\
1.0&0.020&75.719960191&\(2.790\!\times10^{-7}\)&0.104399&\((1,1,1)\)\\
\bottomrule
\end{tabular}
\end{center}
These controls show persistence of admissible stationary representatives, not
regulator independence: the finite-\(a\) energies and profiles visibly depend
on \(\gamma\) and \(\epsilon\).

A deterministic \(2\times10^{-3}\) tangent-plus-scalar perturbation of the
\(N=15,L=6\) initial sample returns to a numerically indistinguishable
same-energy admissible stationary representative: its final energy
differs from the baseline by less than \(10^{-13}\) relatively, has
three-target \(B=(1,1,1)\), and passes the direct stationarity threshold.
This is a local-basin control, not a proof that the reported point is the
global minimizer of its component.

\subsection{Compact-supported smooth-profile scaling}

To avoid a boundary-layer artifact, this regression does not overwrite the
nonzero tail of the canonical AP-E6 profile.  It uses a separate \(C^2\),
degree-one hedgehog with support radius two inside the \(L=6\) box, so the
field already equals the exact vacuum on every boundary face.  Its barrier
energy is:
\begin{center}
\small
\begin{tabular}{rrrr}
\toprule
\(N\)&\(a\)&\(E_{\rm barrier}\)&\(\min q\)\\
\midrule
17&0.375000&30.439452623&0.150863\\
21&0.300000&18.075663943&0.360734\\
25&0.250000&12.263778095&0.551121\\
29&0.214286&8.914172689&0.656982\\
33&0.187500&6.784938174&0.722692\\
\bottomrule
\end{tabular}
\end{center}
A log--log fit gives slope \(2.15910\) with \(R^2=0.999141\).  The four
successive local slopes are \(2.3356,2.1276,2.0695,2.0440\), approaching the
analytic exponent two.  The fit is not applied to the re-relaxed minimizers
and cannot establish their continuum behavior.

\section{Machine gates and changed blocker}

The production card passes \(13/13\) checks.  It writes
\begin{center}
\path{route_f/output/ap_e8_topology_preserving_b1.json},\qquad
\path{route_f/output/ap_e8_topology_preserving_b1.md}.
\end{center}
The following statements are now true:
\begin{enumerate}
\item a topology-preserving finite-grid action is explicitly defined;
\item its finite-energy degree is preserved by theorem;
\item every nonempty fixed-grid \(B=1\) component has an unanchored minimizer;
\item such components and minimizers exist on a sufficiently fine grid
sequence;
\item seven numerical final fields pass a direct unanchored stationarity test,
three-target \(B=1\), and a positive barrier margin.
\end{enumerate}

This changes the AP-E7 conclusion.  The former statement
``no unanchored admissible finite-grid stationary representative has been
found'' is superseded.  The correct new statement is
\begin{equation}
 \boxed{C_{B1}^{\rm finite\ grid}=\mathrm{true},\qquad
 C_{B1}^{\rm physical\ continuum}=\mathrm{false}.}
 \label{eq:changed-blocker}
\end{equation}

\begin{gate}[Continuum and physics remain fail closed]
None of the following has been established: a numerically certified global
minimum, complete line-search-segment admissibility, equicoercivity, Gamma
convergence, a continuum stationary limit, barrier or triangulation
independence, a physical charged-QC2D action, an interacting
gauge--meson--ghost--fermion Hessian, a BRST superdeterminant, four-dimensional
dynamics, HMC, or a quantum continuum limit.
\end{gate}

Consequently the full Lane-B closure condition remains false and no
degree-one Route-E portal is authorized.  The two deferred alternatives---a
complete finite GW/domain-wall target kernel and the actual \(SO(3)\) Yukawa
mapping-torus mod-two index---remain independent research lanes.

\section{Ordered next work}

The next calculations should be ordered by their ability to distinguish a
physical continuum object from a barrier artifact:
\begin{enumerate}
\item Prove an equicoercive estimate or a Gamma-liminf statement for a
controlled scaling \(\gamma(a),\epsilon(a)\), or produce a counterexample
showing barrier-core collapse.  Measure local gradient concentration and a
topological radius, not only total energy.
\item Repeat the stationary family for at least two alternative cube
triangulations and several \(\gamma(a),\epsilon(a)\) trajectories.  Extrapolate
observables jointly in \(a,L,\gamma,\epsilon\); do not extrapolate the barrier
energy as if it were physical.
\item Only after one regulator-stable continuum background is identified,
assemble the full same-action Riemann Hessian including the curvature shift
in Eq.~\eqref{eq:riemann-hessian}, then add interacting gauge/ghost blocks and
the bosonic second variation of one regulated fermion determinant.
\item In parallel, complete either the finite target GW/domain-wall operator
or the actual \(SO(3)\) mapping-torus mod-two index.  A portal remains last.
\end{enumerate}

\section{Reproduction}

From the repository root, run
\begin{verbatim}
python3 route_f/code/scan_ap_e8_topology_preserving_b1.py
\end{verbatim}
The JSON serializer rejects NaN and infinity.  The output records Python,
NumPy, SciPy, source SHA-256 values, the canonical seed checksum, every main
and control case, the three calculus audits, all Boolean theorem and fail-closed
flags, and the \(13/13\) checks.

\bibliographystyle{unsrt}
\bibliography{route_f/tex/ap_e8_topology_preserving_b1}

\end{document}
